% GENERATED FILE. Edit canonical lessons, then run npm run generate:books.
% canonical-source-hash: fnv1a64:a8f0213a0ba84c20
% canonical-lessons: ML-C378-abhayam, ML-C378-abhayartthi, ML-C378-jivahani, ML-C378-palayanam-ceyyuka, ML-C378-tetuka

\chapter{Shelter and Flight}
\label{ch:ml-shelter-and-flight}
\hlchaptermodality{378}

\begin{chapteropening}
\textbf{By the end of this chapter:} \emph{I can say shelter, a refugee, loss of life, to flee and to seek.}

Complete the last lesson of chapter 378: I can say shelter, a refugee, loss of life, to flee and to seek.
\end{chapteropening}

\section[\texorpdfstring{\ml{അഭയം} (abhayaṁ)}{അഭയം (abhayaṁ)}]{\ml{അഭയം} (abhayaṁ) — shelter; refuge}

\label{lesson:ML-C378-abhayam}



\begin{quote}
Before the new one: say the Malayalam for to be cut off, then the Malayalam for to go missing.
\end{quote}

\subsection*{You'll want to know: \ml{അഭയം}}
\textbf{\ml{അഭയം}} — \emph{abhayaṁ} — \textquotedblleft{}shelter; refuge\textquotedblright{}.

\begin{quote}
\ml{വെള്ളപ്പൊക്കത്തിൽ} \ml{ആളുകൾ} \ml{പള്ളിക്കൂടത്തിൽ} \ml{അഭയം} \ml{തേടി}.
\end{quote}

\begin{quote}
\emph{veḷḷappokkattil āḷukaḷ paḷḷikkūṭattil abhayaṁ tēṭi.}
\end{quote}

\begin{quote}
In the flood people took shelter in the school.
\end{quote}

Literally freedom from fear: a safe place in a disaster.

\begin{cousinweb}
\textbf{\ml{അഭയം}} (\emph{abhayaṁ}) is borrowed from Sanskrit \emph{abhaya}, \textquotedblleft{}fearless\textquotedblright{}, and as a noun \textquotedblleft{}safety\textquotedblright{} (CDIAL 507): \emph{a-}, \textquotedblleft{}without\textquotedblright{}, and \emph{bhaya}, \textquotedblleft{}fear\textquotedblright{}.
\end{cousinweb}

\begin{grammarlens}[title={What you've built}]
One more word for telling what happened.
\end{grammarlens}

\subsection*{Your turn}
\begin{itemize}
\raggedright
  \item \emph{Say it:} \emph{abhayaṁ}
  \item \emph{Say it:} \emph{abhayaṁ}, once more
  \item \emph{Say it:} \emph{veḷḷappokkattil āḷukaḷ paḷḷikkūṭattil abhayaṁ tēṭi.} — In the flood people took shelter in the school.
  \item \emph{Recall:} say the Malayalam for to be cut off, then the Malayalam for to go missing, then say \emph{abhayaṁ} again
\end{itemize}

\begin{tcolorbox}[breakable,colback=teal!4,colframe=teal!35!black,title={Before you move on}]
What does \ml{അഭയം} mean? (\textquotedblleft{}Shelter; refuge\textquotedblright{}.) Say it once more.
\end{tcolorbox}
\section[\texorpdfstring{\ml{അഭയാർത്ഥി} (abhayārtthi)}{അഭയാർത്ഥി (abhayārtthi)}]{\ml{അഭയാർത്ഥി} (abhayārtthi) — a refugee}

\label{lesson:ML-C378-abhayartthi}



\begin{quote}
Before the new one: say the Malayalam for to go missing, then the Malayalam for shelter.
\end{quote}

\subsection*{You'll want to know: \ml{അഭയാർത്ഥി}}
\textbf{\ml{അഭയാർത്ഥി}} — \emph{abhayārtthi} — \textquotedblleft{}a refugee\textquotedblright{}.

\begin{quote}
\ml{യുദ്ധം} \ml{കാരണം} \ml{ഒരുപാട്} \ml{അഭയാർത്ഥികൾ} \ml{അതിർത്തി} \ml{കടന്നു}.
\end{quote}

\begin{quote}
\emph{yuddhaṁ kāraṇaṁ orupāṭŭ abhayārtthikaḷ atirtti kaṭannu.}
\end{quote}

\begin{quote}
Because of the war many refugees crossed the border.
\end{quote}

Plural: \ml{അഭയാർത്ഥികൾ}, \emph{abhayārtthikaḷ}.

\ml{അഭയം}, refuge, + \ml{അർത്ഥി}, one who seeks, as in \ml{വിദ്യാർത്ഥി}.

\begin{cousinweb}
\textbf{\ml{അഭയാർത്ഥി}} (\emph{abhayārtthi}) is built from Sanskrit \emph{abhaya}, \textquotedblleft{}safety\textquotedblright{} (CDIAL 507), and \emph{arthin}, \textquotedblleft{}desiring, one who seeks\textquotedblright{} (CDIAL 640), the same ending as in \ml{വിദ്യാർത്ഥി}, \textquotedblleft{}a student\textquotedblright{}.
\end{cousinweb}

\begin{grammarlens}[title={What you've built}]
One more word for telling what happened.
\end{grammarlens}

\subsection*{Your turn}
\begin{itemize}
\raggedright
  \item \emph{Say it:} \emph{abhayārtthi}
  \item \emph{Say it:} \emph{abhayārtthi}, once more
  \item \emph{Say it:} \emph{yuddhaṁ kāraṇaṁ orupāṭŭ abhayārtthikaḷ atirtti kaṭannu.} — Because of the war many refugees crossed the border.
  \item \emph{Recall:} say the Malayalam for to go missing, then the Malayalam for shelter, then say \emph{abhayārtthi} again
\end{itemize}

\begin{tcolorbox}[breakable,colback=teal!4,colframe=teal!35!black,title={Before you move on}]
What does \ml{അഭയാർത്ഥി} mean? (\textquotedblleft{}A refugee\textquotedblright{}.) Say it once more.
\end{tcolorbox}
\section[\texorpdfstring{\ml{ജീവഹാനി} (jīvahāni)}{ജീവഹാനി (jīvahāni)}]{\ml{ജീവഹാനി} (jīvahāni) — loss of life}

\label{lesson:ML-C378-jivahani}



\begin{quote}
Before the new one: say the Malayalam for shelter, then the Malayalam for a refugee.
\end{quote}

\subsection*{You'll want to know: \ml{ജീവഹാനി}}
\textbf{\ml{ജീവഹാനി}} — \emph{jīvahāni} — \textquotedblleft{}loss of life\textquotedblright{}.

\begin{quote}
\ml{ഭാഗ്യത്തിന്} \ml{അപകടത്തിൽ} \ml{ജീവഹാനി} \ml{ഉണ്ടായില്ല}.
\end{quote}

\begin{quote}
\emph{bhāgyattinŭ apakaṭattil jīvahāni uṇṭāyilla.}
\end{quote}

\begin{quote}
Luckily there was no loss of life in the accident.
\end{quote}

A news word; in speech people say \ml{ആരും} \ml{മരിച്ചില്ല}, nobody died.

\begin{cousinweb}
\textbf{\ml{ജീവഹാനി}} (\emph{jīvahāni}) is built from Sanskrit \emph{jīva}, \textquotedblleft{}life\textquotedblright{} (CDIAL 5239), and \emph{hāni}, \textquotedblleft{}abandonment, loss\textquotedblright{} (Monier-Williams).
\end{cousinweb}

\begin{grammarlens}[title={What you've built}]
One more word for telling what happened.
\end{grammarlens}

\subsection*{Your turn}
\begin{itemize}
\raggedright
  \item \emph{Say it:} \emph{jīvahāni}
  \item \emph{Say it:} \emph{jīvahāni}, once more
  \item \emph{Say it:} \emph{bhāgyattinŭ apakaṭattil jīvahāni uṇṭāyilla.} — Luckily there was no loss of life in the accident.
  \item \emph{Recall:} say the Malayalam for shelter, then the Malayalam for a refugee, then say \emph{jīvahāni} again
\end{itemize}

\begin{tcolorbox}[breakable,colback=teal!4,colframe=teal!35!black,title={Before you move on}]
What does \ml{ജീവഹാനി} mean? (\textquotedblleft{}Loss of life\textquotedblright{}.) Say it once more.
\end{tcolorbox}
\section[\texorpdfstring{\ml{പലായനം} \ml{ചെയ്യുക} (palāyanaṁ ceyyuka)}{പലായനം ചെയ്യുക (palāyanaṁ ceyyuka)}]{\ml{പലായനം} \ml{ചെയ്യുക} (palāyanaṁ ceyyuka) — to flee; to take flight}

\label{lesson:ML-C378-palayanam-ceyyuka}



\begin{quote}
Before the new one: say the Malayalam for a refugee, then the Malayalam for loss of life.
\end{quote}

\subsection*{You'll want to know: \ml{പലായനം} \ml{ചെയ്യുക}}
\textbf{\ml{പലായനം} \ml{ചെയ്യുക}} — \emph{palāyanaṁ ceyyuka} — \textquotedblleft{}to flee; to take flight\textquotedblright{}.

\begin{quote}
\ml{യുദ്ധം} \ml{കാരണം} \ml{ആയിരങ്ങൾ} \ml{നാട്ടിൽ} \ml{നിന്ന്} \ml{പലായനം} \ml{ചെയ്തു}.
\end{quote}

\begin{quote}
\emph{yuddhaṁ kāraṇaṁ āyiraṅṅaḷ nāṭṭil ninnŭ palāyanaṁ ceytu.}
\end{quote}

\begin{quote}
Because of the war thousands fled their homeland.
\end{quote}

Past: \ml{പലായനം} \ml{ചെയ്തു}, \emph{palāyanaṁ ceytu}.

The noun \ml{പലായനം} is a flight, an exodus. From where takes -\ml{ഇൽ} \ml{നിന്ന്}.

\begin{cousinweb}
\textbf{\ml{പലായനം} \ml{ചെയ്യുക}} (\emph{palāyanaṁ ceyyuka}) joins \ml{പലായനം}, \textquotedblleft{}flight, running away\textquotedblright{}, borrowed from Sanskrit \emph{palāyana} (Monier-Williams), and \ml{ചെയ്യുക}, \textquotedblleft{}to do\textquotedblright{}.
\end{cousinweb}

\begin{grammarlens}[title={What you've built}]
One more word for telling what happened.
\end{grammarlens}

\subsection*{Your turn}
\begin{itemize}
\raggedright
  \item \emph{Say it:} \emph{palāyanaṁ ceyyuka}
  \item \emph{Say it:} \emph{palāyanaṁ ceyyuka}, once more
  \item \emph{Say it:} \emph{yuddhaṁ kāraṇaṁ āyiraṅṅaḷ nāṭṭil ninnŭ palāyanaṁ ceytu.} — Because of the war thousands fled their homeland.
  \item \emph{Recall:} say the Malayalam for a refugee, then the Malayalam for loss of life, then say \emph{palāyanaṁ ceyyuka} again
\end{itemize}

\begin{tcolorbox}[breakable,colback=teal!4,colframe=teal!35!black,title={Before you move on}]
What does \ml{പലായനം} \ml{ചെയ്യുക} mean? (\textquotedblleft{}To flee; to take flight\textquotedblright{}.) Say it once more.
\end{tcolorbox}
\section[\texorpdfstring{\ml{തേടുക} (tēṭuka)}{തേടുക (tēṭuka)}]{\ml{തേടുക} (tēṭuka) — to seek; to look for}

\label{lesson:ML-C378-tetuka}



\begin{quote}
Before the new one: say the Malayalam for loss of life, then the Malayalam for to flee.
\end{quote}

\subsection*{You'll want to know: \ml{തേടുക}}
\textbf{\ml{തേടുക}} — \emph{tēṭuka} — \textquotedblleft{}to seek; to look for\textquotedblright{}.

\begin{quote}
\ml{ജോലി} \ml{തേടി} \ml{അവൻ} \ml{നഗരത്തിൽ} \ml{പോയി}.
\end{quote}

\begin{quote}
\emph{jōli tēṭi avan nagarattil pōyi.}
\end{quote}

\begin{quote}
He went to the city in search of work.
\end{quote}

Past: \ml{തേടി}, \emph{tēṭi}.

\ml{തേടി} often means in search of. A more literary word than \ml{തിരയുക}.

\begin{cousinweb}
\textbf{\ml{തേടുക}} (\emph{tēṭuka}) is native Dravidian: DEDR 3456 lists Malayalam \emph{tēṭuka}, \textquotedblleft{}to seek, to acquire\textquotedblright{}, beside Tamil \emph{tēṭu}, \textquotedblleft{}to seek, to search for\textquotedblright{}.
\end{cousinweb}

\begin{grammarlens}[title={What you've built}]
One more word for telling what happened.
\end{grammarlens}

\subsection*{Your turn}
\begin{itemize}
\raggedright
  \item \emph{Say it:} \emph{tēṭuka}
  \item \emph{Say it:} \emph{tēṭuka}, once more
  \item \emph{Say it:} \emph{jōli tēṭi avan nagarattil pōyi.} — He went to the city in search of work.
  \item \emph{Recall:} say the Malayalam for loss of life, then the Malayalam for to flee, then say \emph{tēṭuka} again
\end{itemize}

\begin{tcolorbox}[breakable,colback=teal!4,colframe=teal!35!black,title={Before you move on}]
What does \ml{തേടുക} mean? (\textquotedblleft{}To seek; to look for\textquotedblright{}.) Say it once more.
\end{tcolorbox}
